\documentclass[UTF8,zihao=-4]{ctexart}
\usepackage[a4paper,margin=1.65cm]{geometry}
\usepackage{xcolor}
\usepackage[most]{tcolorbox}
\usepackage{enumitem}
\usepackage{hyperref}
\usepackage{fancyhdr}
\usepackage{titlesec}
\usepackage{fontspec}
\usepackage{multicol}
\usepackage{needspace}

\setCJKmainfont{Microsoft YaHei}
\setCJKsansfont{Microsoft YaHei}
\setCJKmonofont{Microsoft YaHei}
\setmainfont{Times New Roman}
\setsansfont{Arial}
\setmonofont{Consolas}

\definecolor{primary}{HTML}{2563EB}
\definecolor{deep}{HTML}{172554}
\definecolor{muted}{HTML}{475569}
\definecolor{lightblue}{HTML}{EFF6FF}
\definecolor{border}{HTML}{BFDBFE}
\definecolor{greenbg}{HTML}{ECFDF5}
\definecolor{green}{HTML}{047857}
\definecolor{warmbg}{HTML}{FFFBEB}
\definecolor{warm}{HTML}{B45309}

\hypersetup{colorlinks=true,linkcolor=primary,urlcolor=primary}
\pagestyle{fancy}
\fancyhf{}
\lhead{大学英语六级写作范文背诵包}
\rhead{10 个高频话题 / 打印版}
\cfoot{\thepage}
\renewcommand{\headrulewidth}{0.4pt}
\setlength{\headheight}{15pt}

\titleformat{\section}{\Large\bfseries\color{deep}}{\thesection}{0.7em}{}[\titlerule]
\titleformat{\subsection}{\normalsize\bfseries\color{primary}}{}{0em}{}
\setlist[itemize]{leftmargin=1.6em,itemsep=0.12em,topsep=0.2em}
\setlist[enumerate]{leftmargin=1.8em,itemsep=0.12em,topsep=0.2em}
\setlength{\parindent}{0pt}
\setlength{\parskip}{0.38em}
\linespread{1.08}

\tcbset{arc=2mm,boxrule=0.5pt,left=2mm,right=2mm,top=1.1mm,bottom=1.1mm}
\newtcolorbox{tipbox}{colback=lightblue,colframe=border}
\newtcolorbox{sentbox}{colback=greenbg,colframe=green}
\newtcolorbox{templatebox}{colback=warmbg,colframe=warm}
\newenvironment{essayblock}{\Needspace{0.46\textheight}\begin{samepage}}{\end{samepage}\vspace{0.2em}}
\newcommand{\essay}[3]{\section{#1}\subsection{#2}\textbf{#3}\par}
\newcommand{\keysentence}[1]{\begin{sentbox}\textbf{可背句：}#1\end{sentbox}}

\begin{document}

\begin{center}
{\Huge \bfseries 大学英语六级写作范文背诵包}\\[0.35em]
{\large 10 个高频话题｜结构化背诵｜打印版}\\[0.3em]
{\small Blog: \href{https://liangqianxing.github.io/posts/cet6-writing-model-essays/}{liangqianxing.github.io/posts/cet6-writing-model-essays/}}
\end{center}

\begin{tipbox}
\textbf{背诵建议：}每篇优先背三类句子：开头立场句、主体转折句、结尾升华句。真正考试时不要逐字硬套，保留句型框架，替换关键词即可迁移到新题目。
\end{tipbox}

\begin{multicols}{2}
\begin{enumerate}
  \item 人工智能与学习
  \item 终身学习
  \item 网络课程
  \item 坚持的重要性
  \item 阅读的价值
  \item 环境保护
  \item 志愿服务
  \item 心理健康
  \item 团队合作
  \item 传统文化
\end{enumerate}
\end{multicols}

\begin{essayblock}
\essay{人工智能与学习}{The Role of Artificial Intelligence in Learning}{In recent years, artificial intelligence has entered almost every corner of education. Some students use AI tools to translate texts, summarize articles, and even polish their essays. In my view, AI can be a powerful assistant, but it should never become a substitute for independent thinking.}

To begin with, AI makes learning more efficient. It can provide instant feedback, explain difficult concepts in different ways, and help students review knowledge according to their own pace. However, excessive dependence on AI may weaken our ability to analyze problems. If students simply copy answers generated by machines, they may finish assignments faster but learn much less.

Therefore, the proper attitude is to use AI wisely. We should ask it for guidance, examples, and inspiration, while making final judgments by ourselves. Only when technology is combined with self-discipline and critical thinking can it truly improve our learning.

\keysentence{AI can be a powerful assistant, but it should never become a substitute for independent thinking.}
\end{essayblock}

\begin{essayblock}
\essay{终身学习}{The Importance of Lifelong Learning}{The world is changing so rapidly that what we learn today may become outdated tomorrow. Under such circumstances, lifelong learning is no longer a choice for a few people, but a necessity for everyone who wants to keep competitive.}

There are several reasons why lifelong learning matters. First, it helps us adapt to social and technological changes. New tools, new industries, and new working methods appear constantly, and those who stop learning may soon fall behind. Second, learning enriches our inner world. Reading, taking courses, and exploring unfamiliar fields can broaden our horizons and make life more meaningful.

As college students, we should develop the habit of continuous learning as early as possible. We can read regularly, learn practical skills online, and reflect on our experiences. In the long run, the ability to learn may be more important than any single piece of knowledge.

\keysentence{Lifelong learning is no longer a choice for a few people, but a necessity for everyone who wants to keep competitive.}
\end{essayblock}

\begin{essayblock}
\essay{网络课程}{Online Courses and Traditional Classrooms}{With the development of the Internet, online courses have become increasingly popular among college students. Some people believe that they will replace traditional classrooms, while others insist that face-to-face teaching is still irreplaceable. In my opinion, the two forms should complement each other.}

Online courses have obvious advantages. Students can study anytime and anywhere, choose resources from top universities, and review difficult parts repeatedly. This flexibility is especially helpful for those who need to arrange their time independently. Nevertheless, traditional classrooms also play an essential role. Direct communication with teachers and classmates can create a better learning atmosphere and encourage students to participate actively.

Therefore, a balanced approach is the best choice. We may use online courses to expand knowledge and use classroom learning to discuss, practice, and receive guidance. When the two are combined properly, students can enjoy both efficiency and interaction.

\keysentence{The two forms should complement each other rather than replace each other.}
\end{essayblock}

\begin{essayblock}
\essay{坚持的重要性}{Persistence Leads to Success}{It is often said that success belongs to those who never give up. This saying is simple, but it contains a truth that has been proved again and again: persistence is one of the most important qualities behind achievement.}

No meaningful goal can be achieved overnight. Whether we are preparing for an exam, learning a language, or doing scientific research, we are likely to meet difficulties and experience failure. At such moments, talent may give us a good start, but persistence determines how far we can go. A person who keeps improving a little every day will finally make remarkable progress.

Of course, persistence does not mean repeating the same mistake blindly. We should also reflect on our methods and adjust our plans when necessary. As long as we keep a clear goal, take steady action, and learn from setbacks, we will be much closer to success.

\keysentence{Talent may give us a good start, but persistence determines how far we can go.}
\end{essayblock}

\begin{essayblock}
\essay{阅读的价值}{The Value of Reading}{In the age of short videos and social media, many people spend less time reading books. However, reading remains one of the most valuable habits we can develop, because it brings benefits that fragmented information can hardly offer.}

First of all, reading deepens our thinking. A good book requires patience and concentration, and it guides us to understand ideas in a systematic way. Besides, reading broadens our vision. Through books, we can communicate with great minds, learn about different cultures, and gain experiences beyond our own life. More importantly, reading can bring inner peace in a noisy world.

To make reading a lasting habit, we do not have to begin with difficult classics. Instead, we can start from books that truly interest us and read a few pages every day. Over time, reading will not only improve our language ability but also shape the way we see the world.

\keysentence{Reading brings benefits that fragmented information can hardly offer.}
\end{essayblock}

\begin{essayblock}
\essay{环境保护}{Protecting the Environment Starts from Small Things}{Environmental protection has become a common concern around the world. Although governments and companies should certainly take major responsibility, ordinary individuals can also make a difference through small daily actions.}

Small actions matter because they can gradually form a powerful social habit. For example, if more people choose public transportation, save electricity, reduce food waste, and avoid unnecessary plastic products, the pressure on the environment will be greatly reduced. In addition, individual behavior can influence others. When one person lives in a greener way, friends and family members may also become more aware of environmental problems.

Therefore, we should not underestimate the value of small things. Environmental protection is not only a slogan, but also a lifestyle. If each of us takes one more responsible step today, society will move closer to a cleaner and more sustainable future.

\keysentence{Environmental protection is not only a slogan, but also a lifestyle.}
\end{essayblock}

\begin{essayblock}
\essay{志愿服务}{The Meaning of Voluntary Work}{Nowadays, more and more college students take part in voluntary work, such as helping in communities, teaching children, and serving at public events. From my perspective, voluntary work is meaningful not only to society but also to volunteers themselves.}

On the one hand, voluntary work helps those in need and makes society warmer. Many social problems cannot be solved by money alone; they also require care, patience, and human connection. On the other hand, volunteers can gain valuable experience from serving others. They learn how to communicate, cooperate, and take responsibility. These abilities are important for both personal growth and future career development.

In short, voluntary work is a bridge between individual growth and social progress. As college students, we should spend some time helping others whenever possible. By giving our time and energy, we may receive a stronger sense of purpose in return.

\keysentence{Voluntary work is a bridge between individual growth and social progress.}
\end{essayblock}

\begin{essayblock}
\essay{心理健康}{Mental Health Deserves More Attention}{In recent years, mental health has attracted increasing attention on campus. Many students face pressure from exams, relationships, future employment, and self-expectations. Therefore, it is necessary for us to take mental health as seriously as physical health.}

Good mental health is the foundation of effective study and a balanced life. When students are anxious or depressed for a long time, they may find it hard to concentrate, communicate, or make decisions. More importantly, ignoring emotional problems may make them worse. Seeking help is not a sign of weakness, but a responsible way to protect oneself.

To improve mental health, students should keep regular routines, exercise properly, and talk with trusted friends, teachers, or counselors when necessary. Universities should also provide accessible psychological services. Only in a supportive environment can students grow with confidence and resilience.

\keysentence{Seeking help is not a sign of weakness, but a responsible way to protect oneself.}
\end{essayblock}

\begin{essayblock}
\essay{团队合作}{The Power of Teamwork}{Teamwork is becoming increasingly important in modern society. Many tasks are too complex to be completed by one person alone, so people with different abilities need to work together toward a shared goal.}

The value of teamwork lies in cooperation and mutual support. In a team, each member can contribute his or her strengths. Some people are good at planning, some are skilled at communication, and others may be creative problem-solvers. When these strengths are combined, the team can often achieve better results than individuals working separately. Moreover, teamwork teaches us to listen, respect different opinions, and handle conflicts in a mature way.

However, effective teamwork does not happen automatically. It requires clear goals, fair division of work, and responsible members. If everyone does his or her part and remains open to others' ideas, teamwork will become a powerful force for success.

\keysentence{When different strengths are combined, the team can often achieve better results than individuals working separately.}
\end{essayblock}

\begin{essayblock}
\essay{传统文化}{Preserving Traditional Culture}{As globalization continues to deepen, traditional culture is facing both opportunities and challenges. Some young people are attracted by foreign cultures and pay less attention to their own traditions. In my opinion, preserving traditional culture is necessary, because it is closely connected with our identity and values.}

Traditional culture records the history, wisdom, and aesthetic taste of a nation. Festivals, literature, music, food, and customs all tell us where we come from and how previous generations understood life. If we lose these cultural roots, society may become richer materially but poorer spiritually. At the same time, preserving tradition does not mean refusing change. We should present traditional culture in modern ways so that more young people can understand and enjoy it.

In conclusion, traditional culture should be protected, inherited, and creatively developed. Only by respecting our cultural roots can we communicate with the world with greater confidence.

\keysentence{Only by respecting our cultural roots can we communicate with the world with greater confidence.}
\end{essayblock}

\section{通用结尾模板}
\begin{templatebox}
\begin{enumerate}
  \item In conclusion, ... is not only important for individuals, but also valuable to society as a whole.
  \item Therefore, we should take practical action instead of treating the issue as an empty slogan.
  \item Only when we combine awareness with action can we make real progress.
  \item As college students, we are expected to develop a responsible attitude and turn it into daily practice.
  \item In the long run, small but continuous efforts will lead to meaningful changes.
\end{enumerate}
\end{templatebox}

\end{document}



